\chapter{The Field Closes}

\section{The second variation}

The machine has become a meter and read the cost of its own act. Before it can ask that meter for the coupling, one more
structure must close. The readings so far have been \emph{first} differences --- how a quantity changes as the machine
turns once. The coupling lives one order deeper, in the \emph{second} difference: how the change itself changes, the
curvature of the machine's cost as it is varied. This section builds that second variation, in the plainest numerical
terms, because it is the object the coupling will be a reading of.

Take a first difference first, to fix the ladder. The machine varies something --- a seed, a parameter of its own
descent --- and reads how a cost responds; that response, the change in cost per change in parameter, is a first
difference, a discrete slope. It is a real reading, but it is not yet the one the coupling needs. Vary the parameter
again and ask how the \emph{slope} responds --- how the first difference itself changes --- and you have a second
difference, the discrete curvature of the cost. Numerical analysis writes it with a three-point stencil, a value at the
centre flanked by values on either side, combined so that the linear part cancels and only the bend remains~\cite{golub2013}. The second variation, $\delta^2$, is this bend: not the slope of the cost but the rate at which the slope
turns, read off a finite stencil the machine can compute.

The construction takes this second difference across \emph{four channels}, and the four-channel shape is the specific
object the coupling reads. The machine does not vary one parameter but a small structured family of them, and it reads
the curvature of its cost across that family --- a four-channel second variation, a curvature with more than one
direction. This four-channel curvature is what a physical gauge, reading the same object, names an electromagnetic
coupling; the construction carries it under that gauge's name, marked here as a reading, not a thing this book builds. In
this gauge it is exactly what the arithmetic says: the second difference of the machine's own cost, taken across four
channels, the discrete curvature of the price the machine pays as its own account is varied. The physical name is handed
off; the curvature is kept.

In the two verbs the second variation is a tange of a tange. The first difference tanges apart how the cost changes; the
second tanges apart how \emph{that change} changes, a distinction drawn one level finer than the first. And the four
channels funge these curvatures into one structured reading --- the four-channel second variation pooled as a single
object, the curvature the coupling will be read from. The machine has told apart not just its cost but the bending of its
cost, and pooled the bendings across four channels into the object the payoff needs. Before that object can be asked for
a number, though, one side of it must be shown present --- which is the next section's insistence.

\section{The magnetic side must be present}

A curvature has more than one side, and the machine may not ask for the coupling until both sides of its four-channel
second variation are shown present. One side the readings so far have made familiar; the other has stayed in the
background, and this section insists it be brought forward. A physical gauge calls the two sides the electric and the
\emph{magnetic}, and names the closing of the second the completion of a field; this gauge reads them as the two halves
of a discrete curvature that must both be nonzero before the curvature is real. The field closes when the second half is
present --- and it must close before the coupling is asked, or the asking would be about half an object.

See why one side alone is not enough. A curvature read in only one direction is like a slope reported without its cross
term --- it captures how the cost bends along one channel and says nothing about how the channels bend against each
other. The four-channel second variation has cross terms, off-diagonal bends that couple one channel to another, and
those cross bends are the second side. A physical gauge reads them as the magnetic part of its field; this gauge reads
them as the off-diagonal entries of the curvature, the mixed second differences that were there all along and had not yet
been shown. To ask the coupling of a curvature with its cross terms unread would be to ask it of a truncated object --- a
half-filled stencil pretending to be full. The magnetic side must be present because the curvature is not closed until
its off-diagonal bends are read.

And here the book's discipline about the continuum reasserts itself, in the one place a physical name might smuggle it
in. A physical gauge's field is a continuous object, filling space; the construction's four-channel curvature is nothing
of the kind. It is a \emph{finite} stencil --- a curvature read across a small, fixed number of channels, at a small,
fixed number of rungs, three of them, the same count-to-three the whole instrument floors at. The physical gauge's
Maxwell-shaped field and its expelled-interior reading are that gauge's continuous names, marked and handed off; what
this book carries is the finite three-rung stencil, a curvature with finitely many entries the machine actually computes.
The field ``closes'' here in the finite sense only: the finite stencil's off-diagonal entries are shown present and
nonzero, the discrete curvature is fully read. Nothing continuous is claimed to close, because the machine holds nothing
continuous --- only the finite stencil, floored at three rungs, marked against the continuum the whole book fences off.

In the two verbs closing the field is a funge the machine completes before it may ask. The second variation tanged the
curvature's bends apart --- diagonal and off-diagonal, the two sides; closing the field funges them into one complete
curvature, the finite stencil filled in both its halves. The machine may not ask the coupling of a half-funged curvature;
it must first pool both sides into the one closed object, the finite three-rung stencil with its cross bends present. The
field closes when that funge closes --- and only then, with the curvature whole and finite, is the machine permitted to
ask. What that asking is, and what it carefully is not, is the chapter's last word.

\section{Asking, not delivering}

The field has closed: the four-channel curvature is whole, its cross terms present, the finite stencil filled. The
machine may now ask for the coupling. This closing section is about the exact character of that asking, because the whole
discipline of the book is concentrated in it. Closing the field \emph{asks} the question; it does not \emph{deliver} the
answer. The chapter that closes the field is not the chapter that hands over the magnitude --- and the difference between
the two is the difference between an honest instrument and a crank.

Be precise about what closing the field accomplishes and what it does not. It accomplishes the \emph{setup}: it assembles
the object the coupling is a reading of --- the whole, finite, four-channel curvature --- and makes it available to be
read. That is a real achievement; without a closed curvature there is nothing to ask the coupling of. But assembling the
object is not reading it. The curvature, closed, poses the question ``what is the coupling?'' the way a fully wired
instrument poses the question ``what does the needle read?'' --- the instrument is ready, the question is well-formed,
and the answer is still to be taken. Closing the field is building the question, completely and honestly. It is not, and
must not be mistaken for, the taking of the reading.

The temptation the section exists to refuse is the temptation to let the setup masquerade as the answer --- to present a
closed field as though the closing itself delivered the magnitude. A crank's instrument does exactly this: it assembles
an impressive apparatus and lets the assembly's completeness stand in for a reading, so that the reader mistakes a
well-built question for an answer. The honest machine keeps them apart. It says: the field is closed, the curvature is
whole, the question is now askable --- and the magnitude is not yet in hand, will be taken by the meter of the last part,
and will come back not as a point but as a bracket floored at the machine's resolution. The closing earns the right to
ask; it earns nothing more; and the machine claims nothing more from it.

This restraint is the whole book's temper, arriving at the threshold of its payoff. Every chapter has leaned toward the
ending and refused to pre-empt it; this one refuses hardest, because it is closest. The reader, watching the field close,
feels the magnitude within reach and wants the chapter to hand it over. The chapter will not. It hands over a closed
curvature and a well-formed question, and it sends the reader into the payoff with the answer still ahead --- because an
answer delivered here, before the meter has read it and bracketed it and floored it, would be a number asserted rather
than measured, exactly the thing the instrument was built not to do. The field closes to ask. The asking is the
chapter's whole gift, and it is a real one: a question, honestly built, held open for the meter to answer.

In the two verbs asking-not-delivering is a tange the machine sets up and declines to complete. Closing the field funges
the curvature into one whole object; asking the coupling would tange a magnitude out of that object --- select, from the
closed curvature, the single number that is its coupling. This chapter performs the funge and stops at the edge of the
tange: it pools the curvature whole and poses the question, but it does not pull the magnitude out, because pulling it out
is the meter's act, taken under the floor, and returned as a bracket. The field closes; the question stands; the tange of
the magnitude is left for the payoff --- where, as every reading in this book has been, it comes back not as a point but
as an interval the machine can repeat.
